\documentclass[14pt]{article}
\usepackage{amsmath,amsthm,amssymb}
\usepackage {mathtext}
\usepackage{amsfonts} 
\usepackage[T1,T2A]{fontenc}
\usepackage[utf8]{inputenc}
\usepackage [english, russian] {babel}
\usepackage[babel=true]{microtype}
\usepackage{caption}
\usepackage{subcaption}
\captionsetup[table]{position=t,justification=raggedleft,slc=off}
\setlength{\abovecaptionskip}{0pt}
\usepackage {geometry}
\usepackage{graphicx}
\usepackage{multirow}
\usepackage{booktabs}
\usepackage{pgfplots}
\usepackage{derivative}
\pgfplotsset{compat=1.18}
\geometry{top=2cm, bottom=2cm, left=2cm, right=1cm}
\title{Отчёт о лабораторной работе №1.3.3. \\ \textbf{чё то про вязкость}}}
\author{Павел Рудачихин, группа Б04-507}
\date{01.04.2026}
\begin{document}
	\maketitle
	\section*{Аннотация}
	
	
	\hspace{4mm} В данной работе изучаются традиционные методы откачки механическим форвакуумным насосом до давления $10^{-2}$ Торр и диффузионным масляным насосом до $10^{-5}$ Торр, а также методы измерения вакуума в этом диапазоне.
	
	
	\textbf{Цель работы:}  1) измерение объёмов форвакуумной и высоковакуумной частей установки; 2) определение скорости откачки системы в стационарном режиме и по ухудшению и улучшению вакуума.
	
	
	\textbf{Оборудование:}  высоковакуумная установка с манометрами: масляным, термопарным и ионизационным.
	
	
	\section*{Теоретические сведения}
	
	\hspace{4mm} Производительность насоса определяется скоростью откачки $W$ (л/с). Разделим вакуумную систему на две части: откачиваемую и насос (к которому отнесём, помимо непостредственно насоса, трубопроводы и краны). Пусть $Q_д$ - количество газа, десорбирующееся с поверхности откачиваемого объёма в единицу времени, $Q_и$ - то же через течи, а $Q_н$ - поток газа, поступающего из насоса обратно в систему. Основное уравнение откачки имеет вид
	
	\begin{equation}
		-VdP = (PW - Q_д - Q_н - Q_и)dt
	\end{equation}
	
	где $V$ - откачиваемый объём. При достижение предельного вакуума
	
	\begin{equation}
		\odv{P}{t} = 0 \Rightarrow P_{пр}W = Q_д + Q_н + Q_и
	\end{equation}
	
	Выразим скорость откачки для предельного вакуума
	
	\begin{equation}
		W = \frac{\sum Q_i}{P_{пр}}
	\end{equation} 
	
	Обычно все Q можно считать постоянными.
	
	Проинтегрировав уравнение (1), считая скорость откачки постоянной, получаем:
	
	\begin{equation}
		P - P_{пр} = (P_0 - P_{пр}) \exp \left( - \frac{W}{V}t \right)
	\end{equation}
	
	где $P_0$ - начальное давление, которое велико по сравнению с предельным, поэтому 
	
	\begin{equation}
		P - P_{пр}= P_0 \exp \left( - \frac{W}{V}t \right)
	\end{equation}
	
	Введём постоянную времени откачки $\tau := V/W$, которая является мерой эффективности системы.
	
	Рассмотрим подробнее насосную часть, которая состоит, помимо насоса, из трубок и кранов. Каждый элемент $i$ (насос или трубка) характеризуется своей пропускной способностью $C_i$. Она аналогична проводимости в электрических цепях, и при последовательном соединении элементов имеем:
	
	\begin{equation}
		\frac{1}{W} = \frac{1}{W_н} + \sum_{i=1}^n \frac{1}{C_i}
	\end{equation}
	
	Для количества газа, протекающего через трубу в условиях высокого вакуума справедлива формула:
	
	\begin{equation}
		\odv{PV}{t} = \frac{4}{3} r^3 \sqrt{\frac{2 \pi RT}{\mu}} \frac{P_2 - P_1}{L}
	\end{equation}
	
	Пренебрегая $P_1$ найдём пропускную способность трубы
	
	\begin{equation}
		C = \frac{4}{3} \frac{r^3}{L} \sqrt{\frac{2 \pi RT}{\mu}} 
	\end{equation}
	
	
	
	\section*{Экспериментальная установка}
	
	\begin{figure}[h!]
		\centering
		\includegraphics[width=0.6\linewidth]{setup}
		\caption{Схема установки:
			ВБ - высоковакуумный баллон;
			ФБ - форвакуумный баллон;
			ВН - высоковакуумный диффузионный насос;
			М - масляный манометр;
			И - ионизационный манометр;
			$M_1$, $M_2$- термопарные манометры;
			ФН - форвакуумный насос;
			К - краны.
		}
		\label{fig:setup}
	\end{figure}
	
	
	\section*{Результаты измерений}
	
	\subsection*{Определение объёма частей установки}
	\hspace{4mm} Форвакуумным насосом установка откачана до $0,019$ Торр.  В кранах К5 и К6 и соединяющем их капилляре заперто $V_0 = 50 \ см^3$ воздуха при атмосферном давлении $P_A = 99,67$ кПа. Последовательно открыв краны К5 и К3 заполним этим воздухом сначала форвакуумную часть $V_ф$, затем высоковакуумную $V_в$. При этом измерим уровни масла в манометре $H$ и $h$ соответственно (см. табл. 1).
	
	\begin{table}[h!]
		\centering
		\captionbox{Измерение давления масляным манометром}{
			\begin{tabular}{|r|r|r|r|r|r|r|}
				\hline
				№        & $H_1$, см       & $H_2$, см       & $\Delta H$, см      & $h_1$, см      & $h_2$, см       & $\Delta h$, см \\
				\hline
				1        & 38,3     & 11,2     & 27,1     & 33,9     & 16,4     & 17,5 \\
				2        & 38,7     & 11,6     & 27,1     & 34,1     & 16,6     & 17,5 \\
				\hline
		\end{tabular}}
		\label{dh}%
	\end{table}%
	
	Зная плотность масла $\rho = 885 \ кг/м^3$, найдём разницу давлений $p = \rho g \Delta h$, равное давлению внутри установки. 
	
	По закону Бойля-Мариотта определим объёмы частей установки:
	
	\begin{equation}
		PV = const \Rightarrow V_ф = \frac{P_A}{\rho g \Delta H} V_0; \ \ V_в = V_{полн} - V_ф = \left(\frac{1}{ \Delta h} - \frac{1}{\Delta H}\right) \frac{P_A V_0}{\rho g}
	\end{equation}
	
	Получим
	
	\begin{equation}
		V_ф = (2120 \pm 15) \ мл \ \ \ V_в = (1160 \pm 50) \ мл 
	\end{equation}
	
	\subsection*{Получение высокого вакуума и измерение скорости откачки}
	
	\begin{equation*}
		P_{пр} = 9 \cdot 10^{-5}  \ мм. \ рт. \ ст.
	\end{equation*}
	
	Включением и отключением откачки с помощью крана К3 по улучшению и ухудшению вакуума соответственно можно определить скорость откачки и $Q_н$. 
	
	График зависимости давления от времени в течение эксперимента представлен на рис. 2.
	
	\begin{figure}[h!]
		\centering
		\includegraphics[width=0.6\linewidth]{o}
		\caption{График зависимости P(t)}
		\label{fig:o}
	\end{figure}
	
	
	Спрямим график в координатах $\ln P (t)$ на участке улучшения вакуума (см. рис. 3).
	
	\begin{figure}[h!]
		\centering
		\includegraphics[width=0.49\linewidth]{1down}
		\includegraphics[width=0.49\linewidth]{2down}
		\caption{График зависимости $\ln (P - P_{пр}) (t)$ для улучшения вакуума}
		\label{fig:o2}
	\end{figure}
	
	Далее по методу наименьших квадратов определим угловой коэффициент $k$ и свободный член $b$ для наилучшей прямой:
	
	\begin{equation*}
		k = (-0,095 \pm 0,004) \ c^{-1} \ \ \ b = (2,39 \pm 0,03)
	\end{equation*}
	
	Логарифмируя (5) получаем
	
	\begin{equation*}
		k = -\frac{W}{V} \ \ \ b = \ln P_0
	\end{equation*}
	
	\begin{equation*}
		W = -kV = \boxed{(140 \pm 10) \ мл/с}
	\end{equation*}
	
	При ухудшении вакуума уравнение (1) преобразуется так
	
	\begin{equation}
		\odv{P}{t} = \frac{Q_д + Q_и}{V}
	\end{equation}
	
	Аналогично определяя угловой коэффициент прямой на графике $P(t)$ для ухудшения вакуума (см. рис. 4), найдём значение этой производной
	
	\begin{equation*}
		\odv{P}{t} = (99 \pm 2) \cdot 10^{-5} \ Па/с
	\end{equation*}
	
	\begin{figure}[h!]
		\centering
		\includegraphics[width=0.49\linewidth]{1up}
		\includegraphics[width=0.49\linewidth]{2up}
		\caption{График зависимости P(t) для ухудшения вакуума}
		\label{fig:o3}
	\end{figure}
	
	По формуле (2) оценим $Q_н$:
	
	\begin{equation*}
		Q_н = P_{пр} W - V\odv{P}{t} = \boxed{(0,15 \pm 0,05) \cdot 10^{-5} \ Па\cdot м^3/c}
	\end{equation*}
	
	Открыв кран К6 введём искусственную течь в установку. Установившееся давление 
	
	\begin{equation*}
		P_{уст} = 1,6 \cdot 10^{-4}  \ мм. \ рт. \ ст.
	\end{equation*}
	
	Давление со стороны форвакуумного насоса
	\begin{equation*}
		P_{ф} = 1,4 \cdot 10^{-3}  \ мм. \ рт. \ ст.
	\end{equation*}
	
	Когда капилляр перекрыт и открыт соответственно
	
	\begin{equation}
		P_{пр} W = Q, \ \ \ P_{уст} = Q + \odv{PV}{t} \Rightarrow P_{уст} = P_{пр} W + \odv{PV}{t}
	\end{equation}
	
	По формуле (7) и (8) получим
	
	\begin{equation}
		P_{пр} W = Q, \ \ \ P_{уст}W = Q + \odv{PV}{t} \Rightarrow P_{уст}W = P_{пр} W + (P_ф - P_{уст}) C \Rightarrow W = С \frac{P_ф - P_{уст}}{P_{уст} - P_{пр}} = \boxed{(10 \pm 7) \ мл/с}
	\end{equation}
	
	где производительность капилляра
	
	\begin{equation*}
		C = (5,8 \pm 1,7) \cdot 10^{-7} \ м^3/с
	\end{equation*}
	
	
	\section*{Оценка погрешностей}
	\subsection*{Инструментальные погрешности}
	\hspace{4mm} Оценим погрешность электронного барометра $\sigma(P_A) = 0,03$ кПа.
	
	Значения плотности масла и объёма запертого между К5 и К6 воздуха приведены без погрешностей, поэтому будем считать их измеренными с высокой точностью и пренебрежём их погрешностью.
	
	Шкалой масляного манометра является стальная линейка, её погрешность $\sigma_h = 1$ мм.
	
	
	
	\subsection*{Объёмы частей установки}
	\hspace{4mm} По формуле 4 оценим относительную погрешность измерения объёма:
	
	\begin{equation}
		\varepsilon_V = \sqrt{\varepsilon_{P_A}^2 + \varepsilon_{\Delta h}^2}
	\end{equation}	
	
	где $\varepsilon_{\Delta h} = 2\sigma_h / \Delta h$
	
	Относительная погрешность объёма форвакуумной части составила 0,7\%, что гораздо меньше высоковакуумной части - 4\%. Это обусловлено тем, что объём высоковакуумной части определяется по разности полного объёма и форвакуумной части, причём погрешность полного объёма больше, потом что разность высот уровня масла меньше. 
	
	\subsection*{Скорость откачки по улучшению и ухудшению вакуума}
	
	\hspace{4mm} Инструментальная погрешность измерения давления гораздо меньше случайной погрешности, которая определена по методу наименьших квадратов.
	
	\subsection*{Скорость откачки по искусственной течи}
	
	\hspace{4mm}По формуле (8)
	
	\begin{equation}
		\varepsilon_C = \sqrt{9\varepsilon_r^2 + \varepsilon_L^2}
	\end{equation}
	
	По формуле (13)
	
	\begin{equation}
		\varepsilon_W = \sqrt{\varepsilon_C^2 + \varepsilon_1^2 +  \varepsilon_2^2}
	\end{equation}
	
	где $\varepsilon_1$ и $\varepsilon_2$ относительные погрешности значений $P_ф - P_{уст}$ и $P_{уст} - P_{пр}$ соответственно
	
	
	Отметим, что относительная погрешность измерения радиуса велика, притом входит в формулу с множителем 9 и вносит наибольший вклад в погрешность измерения W методом искусственной течи.
	
	
	
	
	\section*{Заключение}
	
	\hspace{4mm} Цели работы достигнуты: измерены объёмы высоковакуумной и форвакуумной частей установки, измерена скорость откачки диффузионного насоса двумя методами.
	
	Метод измерения скорости откачки с помощью искусственной течи оказался неточен (относительная погрешность почти 100 \%). Основным источником погрешности стало измерение диаметра капилляра и разности давлений. 
	
	Метод измерения скорости откачки по ухудшению и улучшению вакуума относительно точен (погрешность меньше 10 \%) и нагляден, но трудоёмок из-за сложности оцифровки видеозаписи показаний цифрового манометра.
	
	
\end{document}